\chapter{Conclusion}
\label{chap:conclusion}

% ============ GUIDELINE AUDIT 2026-07-25 -- CONCLUSION ============
% Feeds the 15% "Testing, evaluation, critical analysis & conclusions" criterion and the examiner's
% free-text boxes. Full audit: docs/TCD_Dissertation/GUIDELINE_COMPLIANCE.md
%
% The Report Guidelines are specific about this chapter: "The final chapter should be divided into
% two parts: namely, conclusions and future work. The conclusions should be a CRITIQUE of your work:
% it should examine what has been achieved in the project and relate this to the initial aim of the
% project. ANY PROBLEMS ENCOUNTERED IN THE COURSE OF THE PROJECT SHOULD BE MENTIONED HERE."
%
% Future Work is thorough and well-motivated -- six subsections, each tied to a stated limitation.
% Section~\ref{sec:conclusion-answer} relates the result back to the initial aim. Three gaps:
%
% (1) MISSING: THE PROBLEMS-ENCOUNTERED HALF. This is the clearest outright guideline gap in the
%     dissertation. The marking sheet gives the examiner a box headed "Challenges: What were the
%     major difficulties encountered (technical or other)? Comment on the appropriateness of the
%     solutions, if any, offered. Mention any difficulties you are aware of that may have adversely
%     affected the research... Include both difficulties beyond the student's control such as
%     equipment unavailability, and other factors in the student's work."
%     The examiner must fill that box and the dissertation currently gives them nothing to fill it
%     from. Limitations (Chapter~\ref{chap:discussion}) answer a different question -- what the work
%     cannot claim -- not what went wrong and how it was handled.
%     ACTION: add \section{Challenges Encountered} before \section{Future Work}. A draft grounded
%     only in difficulties already evidenced elsewhere in this dissertation is given below; verify
%     each against what actually happened, cut anything that overstates, and add the difficulties
%     only the author knows about. Note this section HELPS rather than hurts: every difficulty below
%     is paired with the design response it produced, which is exactly what "comment on the
%     appropriateness of the solutions offered" invites.
%
% (2) NO FIGURE OR TABLE IN THIS CHAPTER. tab:hypothesis-verdicts is already drafted directly below
%     and commented out. It is the standard "objectives revisited" device and it maps one-to-one
%     onto the 70+ descriptor "clear link with ALL chapters in keeping with research
%     questions/objectives". ACTION: uncomment it and add the pointer sentence its comment supplies.
%     If it goes in, apply the [COMPRESS 2026-07-24] recommendation below as well, since the
%     per-hypothesis prose then duplicates the table.
%
% (3) ORPHANED ETHICS APPENDIX. Appendix~\ref{chap:ethics-statement} exists and is never referenced
%     from the body -- the one reference was swept into a comment block at
%     Section~\ref{subsec:conclusion-future-user-studies} (the [FLAG] there already records this).
%     An unreferenced ethics statement reads as an oversight rather than a decision. ACTION:
%     reinstate the reference, and consider a second pointer from Section~\ref{sec:limitations}.
%
% Also: citation misattribution site 4 of 4 is in this chapter, flagged inline at
% Section~\ref{subsec:conclusion-future-grpo}; and four live contractions, flagged inline.
%
% ---------------------------------------------------------------------------------------------
% [DRAFT 2026-07-25 | testing (15%) + examiner's Challenges box] PROPOSED NEW SECTION.
% Place immediately before \section{Future Work}. Every difficulty named is already evidenced
% elsewhere in this dissertation, so nothing here is invented -- but VERIFY each against what
% actually happened before making it live, and add anything missing that only you know about
% (hardware failures, lost runs, API outages, time lost to a dead end, supervisor redirections).
% Register: third person, matching the rest of the document. To apply, delete the "% " prefixes.
%
% \section{Challenges Encountered}
% \label{sec:conclusion-challenges}
%
% Several difficulties shaped the course of this work, and the design decisions they forced are as
% much a part of the result as the measurements themselves.
%
% The first was a defect in the training data that was not visible until after the models had been
% trained. During trajectory construction the teacher's reasoning was dropped on the turns that
% follow a tool result and back-filled with a short template, so a substantial share of the training
% targets carry no genuine externalised reasoning at all. This was diagnosed only when the
% structural diagnostics of Chapter~\ref{chap:experiments-results} showed the empty-reasoning rate
% rising to 91.3\%, at which point regenerating the corpus and retraining every condition was beyond
% the remaining time budget. The defect is therefore reported rather than repaired, its consequences
% are separated from those of model scale as far as the evidence allows
% (Section~\ref{subsec:h4-failure-boundaries}), and the repair is specified as the first controlled
% experiment future work should run (Section~\ref{subsec:conclusion-future-data}). It is the single
% largest threat to the interpretation of Experiment~3, and it is named as such rather than
% minimised.
%
% The second was the compute budget. All training and generation ran on a single rented
% consumer-class GPU, which closed off several options: full fine-tuning was replaced by sixteen-bit
% LoRA, the five conditions had to be benchmarked one after another rather than in parallel, and
% every run had to survive disconnection. The response was to make the pipeline restartable rather
% than fast, streaming each completed benchmark item to an append-only file so that an interrupted
% run loses nothing, and decoupling judging from generation entirely so that the GPU could be
% released before any scoring began. That decoupling began as a cost measure and became a
% methodological asset, since it allows a saved report to be re-graded by a different judge without
% re-running a model.
%
% The third was that the first measurement instrument proved inadequate on its own. The deterministic
% probes are exactly repeatable, but they under-credit a correct answer phrased in unexpected words,
% and the absolute judge grades every answer against a constitutional ideal that a model of this size
% rarely reaches, so its scores compress towards the bottom of the scale for all five conditions at
% once. Neither instrument alone could distinguish the conditions in the way a reader comparing the
% transcripts could. The response was the dual-lens design of Section~\ref{subsec:mr-judge-validity}:
% a blind comparative judge was added to answer the reader's question directly, with the rule score
% retained beside every judge-based number as a judge-free anchor. The disagreement between the two
% lenses then turned out to carry the third hypothesis's finding
% (Section~\ref{subsec:h3-prompt-engineering}), so an instrument problem became a result.
%
% The fourth was that reinforcement learning could not be attempted. Post-training by reinforcement
% learning is unstable below one billion parameters, and its compute and time requirements were in
% any case beyond what was available, so the scope was fixed to supervised fine-tuning in agreement
% with the supervisor (Section~\ref{sec:scope-and-assumptions}). This bounds what the dissertation
% can say: supervised fine-tuning can only imitate behaviour its examples already demonstrate, so the
% question of whether rewarding outcomes would sidestep the safety tax or merely relocate it remains
% open, and is set out as future work in Section~\ref{subsec:conclusion-future-grpo}.
%
% The fifth was that generating the training data against live tools proved brittle. Because the
% teacher's tool calls were executed for real rather than simulated, each trajectory depended on a
% third-party search service and a sandboxed execution environment behaving predictably, and neither
% always did. Rather than removing the dependency, the failure modes were folded into the corpus as
% deliberate adversarial environments, so that the model was trained on tool timeouts and missing
% tools as well as on successful calls. The alternative, generating trajectories against simulated
% tool results, would have been more reliable but would have taught the model to expect a world that
% does not exist.
% ---------------------------------------------------------------------------------------------
% ==================================================================

The dissertation asked whether a written constitution can be taught to a model small enough to run on a user's own device, so that trustworthy and personalised help need not be bought at the price of sending personal data to the cloud (Figure~\ref{fig:architecture}). This closing chapter gathers what the evidence supports, states plainly what it does not, and marks out the work that would carry the framework further.

\section{Summary of Contributions}
\label{sec:conclusion-contributions}

% [CUT 2026-07-24 | point 1] ACTION: CONSIDER (judgement call, not mechanical): compress the first contribution -- the live paragraph starting "The dissertation provides five contributions" -- to roughly one sentence and let the Discussion carry the argument.
% (Context: the "model-agnostic framework, not a verdict" message duplicates Discussion Section~\ref{subsec:transferable-framework} (its fullest statement) and Introduction Section~\ref{sec:contributions}. The Discussion AUDIT block preceding Section~\ref{subsec:transferable-framework} already tracks this five-way repetition and recommends shortening this contribution to a single asserting clause; a contributions list should assert the contribution, not re-argue it.)
The dissertation makes five contributions. The first, and the primary one, is a reusable, \emph{model-agnostic framework} for measuring and improving constitutional compliance in a small language model, rather than a production-ready model in itself. It comprises a set of deterministic probe suites that test behaviour identically on every run. A yes-or-no verdict on whether the particular Qwen3-0.6B model used here is trustworthy is superseded the moment a more capable small model is released; the procedure for reaching that verdict is not, and it can be re-run on the next model to locate that model's boundary.

On top of that method sit three academic contributions. The first is a systematic constitutional evaluation below one billion parameters: the state-of-the-art survey (Chapter~\ref{chap:state-of-the-art}) placed every prior constitutional study at eight billion parameters and above, and here the language-model judge is treated as a proper measurement instrument whose rubrics, anchors and bias controls are all stated openly. The second is a measurement of the \emph{safety tax} at this scale --- a specific case of the broader \emph{alignment tax} and the central finding of the work --- showing that constitutional fine-tuning buys more compliant behaviour at a measurable cost to the model's visible reasoning, so that it redistributes capability rather than adding it (Chapter~\ref{chap:experiments-results}). The third is a reproducible benchmark that can be run on any Hugging Face-compatible model.

The fifth contribution is architectural. Where the existing literature pairs a small planner with a large, cloud-hosted executor (Section~\ref{sec:multi-model-and-thinker-executor-architectures}), this work instead designs and evaluates a \emph{Thinker--Executor} architecture in which neither half is a large model: both run on the device, and they run one after the other, so that only a single small model is ever computing at any instant and the dependency on a cloud executor is removed, preserving the on-device privacy guarantee the whole dissertation is built around. Across all five, the framework is the contribution the dissertation most rests on, and the verdict on any trained model is just one of the quantities it measures.

\section{Answering the Research Question}
\label{sec:conclusion-answer}

The research question asked in the dissertation whether a framework can teach a small language model the principles of a written constitution by training them into its weights, wrap that model in a deterministic serving layer, and so deliver trustworthy, personalised and privacy-preserving assistance while running entirely on the user's own device. The answer is a \emph{qualified} yes, and the qualification carries the finding, because the boundary between what is feasible at this scale and what is not is itself the result. It is answered hypothesis by hypothesis.

% [DRAFT 2026-07-18 -- at-a-glance verdict table closing the research-question loop (an examiner looks
% for a clear link back to the stated hypotheses; this is the "objectives revisited" device). Numbers from
% Table~\ref{tab:experiment-ladder} and the head-to-head leaderboard. Uncomment and add a pointer such as
% "Table~\ref{tab:hypothesis-verdicts} summarises the verdicts" to the sentence above. The prose that
% follows then unpacks each row.]
% \begin{table}[H]
%   \centering\small
%   \caption{Verdict on each hypothesis; the evidence column cites the effect that decides it, developed in
%     full in Chapter~\ref{chap:discussion}.}
%   \label{tab:hypothesis-verdicts}
%   \begin{tabular}{@{}p{3.5cm}p{2.7cm}p{6.2cm}@{}}
%     \toprule
%     Hypothesis & Verdict & Deciding evidence \\
%     \midrule
%     First: distillation improves compliance & Supported (qualified) & Combined constitution score $0.356 \to 0.485$, the largest single effect in the study; the gain is uneven across principle families. \\
%     \addlinespace
%     Second: the gain costs reasoning & Supported & Empty-reasoning rate $1.4\% \to 91.3\%$; mean trace length $1{,}309 \to 150$ characters. \\
%     \addlinespace
%     Third: prompt engineering with tools underperforms fine-tuning & Mixed (lens-dependent) & Absolute lens favours the fine-tuned model; the blind comparative lens ties the two ($0.589$ vs $0.583$), with the tool baseline leading the trust-critical Tier-1 outcomes. \\
%     \addlinespace
%     Fourth: the failure boundary is systematic & Supported & Failures cluster by principle family rather than by question; the reasoning-intensive principles fail together and in the same direction. \\
%     \addlinespace
%     Fifth: splitting thinking from execution restores capacity on-device & Not supported on its headline claim & Reasoning partly restored (empty-trace $91.3\% \to 36.2\%$; the only condition that clarifies) but combined score $0.367 < 0.485$ and the weakest long-conversation stability ($0.100$). \\
%     \bottomrule
%   \end{tabular}
% \end{table}

The first two hypotheses turn out to be two faces of one outcome. Constitutional distillation, which trained the model on worked examples until the constitution was encoded in its weights, did measurably improve compliance over the untouched baseline, and the staged design shows that this gain belongs to the content of the training data rather than to the act of fine-tuning itself (first hypothesis); but it did so at a price, because the very same training visibly displaced the model's externalised reasoning (second hypothesis), so that the model no longer showed its thinking before answering. This is the \emph{safety tax}: at this scale a gain in one behaviour is paid for by a loss in another, so the improvement is a redistribution and not a free addition. The third hypothesis then asked whether prompt engineering with tool access can match constitutional fine-tuning, and the answer is layered: under the absolute lens the fine-tuned model stays ahead, while the blind comparative lens puts the two at parity but with different strengths, grounding winning the trust-critical outcomes and training winning the reasoning substrate (third hypothesis).

The last two hypotheses turn the harder findings into the actual contribution. Where the fine-tuned model underperformed, it did so \emph{systematically} rather than at random, in a way one could predict in advance from the kind of task and the principle being probed (fourth hypothesis), and that is exactly what makes a located failure a result rather than noise. And by dividing the work across a Thinker that reasons and an Executor that acts, the architecture recovered precisely the visible reasoning the single model had given up, together with the ask-before-assuming behaviour, though it did not become compliant enough overall, paying a measured cost in long-conversation stability (fifth hypothesis). Pulling this together, constitutional AI below one billion parameters is feasible for the structural, safety-oriented and tool-discipline behaviours, sitting layered between grounding and training on the trust-critical outcomes, and it is not yet feasible for the reasoning-intensive principles; and locating that boundary precisely, rather than declaring a verdict for or against, is exactly what this dissertation set out to do.

\section{Future Work}
\label{sec:conclusion-future}

% [STYLE 2026-07-24 | DROP-IN] ACTION: two REPLACEs -- (a) in the GRPO subsection REPLACE "That makes it a open enhancement for future to set out rewards for the models correct behaviour." -> "It is therefore an open extension: rewarding correct behaviour directly, rather than only imitating it."; (b) in the validate-and-retry subsection, apply the replacement quoted below.
% [STYLE 2026-07-24 -- two Future-Work subsections have broken sentences; suggested professional versions:
%  (a) GRPO subsec (Section~\ref{subsec:conclusion-future-grpo}): replace "That makes it a open enhancement
%      for future to set out rewards for the models correct behaviour." with -- "It is therefore an open
%      extension: rewarding correct behaviour directly, rather than only imitating it." (See also the
%      RAGEN/Rainone re-citation flagged in the Introduction AUDIT-FIX, which applies to this subsection.)
%  (b) Validate-and-retry subsec (Section~\ref{subsec:conclusion-future-validate-retry}): replace "provides
%      better self assment before model adjust the weights internally, however since a 0.6B model variant
%      cannot reliably critique itself a better external mechanism to validate-and-retry is a must." with --
%      "assesses a candidate answer before the weights are updated; because a 0.6B model cannot reliably
%      critique itself, that assessment is better supplied externally, by a deterministic validate-and-retry
%      layer."]
Each of the next steps falls directly out of a limitation or scope exclusion above, ordered roughly from the cheapest test of the framework's generality to the most open-ended.

\subsection{Extending to Other Small Models}
\label{subsec:conclusion-future-extensions}

A framework that calls itself model-agnostic has only made a claim until it is run on a second model, so the most direct and cheapest next step is to take the probe suites exactly as they stand and re-run them, without changing the benchmark itself, on a handful of other small models, for instance Phi-3.5-mini, Gemma-2B and SmolLM2-1.7B. These are chosen to vary two things at once, the parameter count and the architectural family, so that the results can begin to tell the two apart.

What this buys is precisely what a single model cannot give. The safety-tax boundary located here is, for now, one point on one model; re-running the framework would show whether that boundary tracks raw \emph{size} (does a slightly larger model push it back in a smooth, predictable way?) or \emph{architecture} (does a differently-built model of similar size sit somewhere else entirely?). The framework was deliberately built to be re-run, so this is a low-cost first test of its generality and a direct response to the single-model limitation conceded in the discussion; the answer it returns is the first real evidence that the method, and not just one set of numbers, travels.

\subsection{Improving Training Data Quality}
\label{subsec:conclusion-future-data}

% Re write this in a way dataset was a synthetic generation putting in early biasness and caveats, a better human judged data would be a better choice here
The obvious follow-up is ``more and better data'', but the measurements allow a good deal more precision than that, because the displacement found is partly a \emph{data-capture defect} rather than a simple shortage of examples. A large share of the single-model training targets carry no genuine externalised reasoning at all: on the turns that follow a tool result, the teacher's reasoning was dropped during trajectory construction and back-filled with a short template, so the model was shown an answer with the thinking already removed. A model can only imitate what its examples actually contain, so where the reasoning was absent from the data, the model learned to omit it.

% [DRAFT -- rewrites the paragraph above per your note (synthetic generation introduced bias and
% caveats early; human-judged data would be the better choice). House register, not the casual [STYLE]
% one. Swap in / merge if approved.]
% A limitation to confront before any refinement is that the training set was generated synthetically
% by the pipeline itself, so the generator's own biases and shortcuts were built into the data from the
% outset rather than introduced later. The clearest of these is a data-capture defect: on the turns that
% follow a tool result, the teacher's reasoning was dropped during trajectory construction and replaced
% by a short template, so a large share of the targets carry no genuine externalised reasoning, and a
% model can only imitate what its examples contain. The stronger remedy, and the better choice where
% resources allow, is to move away from a purely synthetic corpus towards human-judged data: a curated
% set whose reasoning and answers are written or vetted by people removes the generator's inherited
% biases at the source, at the cost of the time and expense that made the synthetic route attractive in
% the first place. A cheaper automated step in the same direction is reasoning-preserving
% self-distillation, in the manner of the on-policy self-distillation that Fu et al.\ show reduces the
% safety tax on small models \cite{fu2026reducingsafetytaxllm}, which regenerates the reasoning targets
% from the base model's own native reasoning rather than a template. Either way, the controlled test the
% work leaves open is the same: if restoring faithful reasoning closes the gap, the shortfall was a data
% problem; if it survives, it was scale (Chapter~\ref{chap:discussion}).

The concrete, cheap fix follows directly from that diagnosis, and it is not ``more data'' but better-captured data. Reasoning-preserving self-distillation, in the manner of the on-policy self-distillation that Fu et al.\ show reduces the safety tax on small models \cite{fu2026reducingsafetytaxllm}, would regenerate the \texttt{<think>} targets from the base model's own native reasoning, run locally with no paid teacher, and graft the constitution-compliant answer onto that genuine reasoning rather than onto a template. Done this way the next experiment is a clean adjudication of the question the work leaves open: if restoring the captured reasoning closes the gap, the displacement was a data problem after all; if the collapse survives even faithfully-captured data, it was scale, exactly as the discussion argues (Chapter~\ref{chap:discussion}).

Two sharpenings follow from the comparative-lens results. First, a coverage audit should be the first step of the loop: beyond capture fidelity, the comparative judge's evidence suggests some probe behaviours have little counterpart in the training corpus at all, precise multi-step computation (calculus-style and chained-arithmetic questions) being the clearest candidate. Before generating new data, the probe taxonomy should be enumerated against the corpus composition, the coverage table published, and the uncovered cells left to drive generation, which turns ``more data'' into a measurable checklist and connects to the fourth hypothesis's boundary discussion. Second, refinement, not only regeneration: the reasoning which was captured should also be refined, since the judge repeatedly graded answers as right direction, failed execution (grade C), that is, the behaviour was chosen correctly and the carry-through failed. A refinement pass that keeps the behavioural choice and tightens the execution steps targets exactly that grade band.

\subsection{A Constitutional Validate-and-Retry Layer}
\label{subsec:conclusion-future-validate-retry}

Anthropic's critique-and-revise loop \cite{bai2022constitutionalaiharmlessnessai} assesses a candidate answer before the weights are updated; because a 0.6B model cannot reliably critique itself, that assessment is better supplied externally, by a deterministic validate-and-retry layer. Concretely it would validate a candidate answer against the deterministic probes, inject a short corrective prompt naming the broken principle, and regenerate under a bounded retry cap, with a persistent scratchpad carrying reasoning state across turns and an explicit read-and-write step against the user memory.

The decisive design requirement is that every one of these components be switchable, so that the same weights can be served once with the layer on and once with it off, everything else held fixed. That on-versus-off comparison over identical weights is the single most direct way to measure whether a runtime layer can recover displaced capability that retraining cannot.

\subsection{User Studies and Judge Validation}
\label{subsec:conclusion-future-user-studies}

User studies and evaluation are deliberately left out. Every result here is a statement about how the model \emph{behaves}, measured by probes and an LLM judge; however, none of it is a statement about how a person \emph{perceives} the system. The probes can show that the model behaves trustworthily mathematically, but only a user can say whether it feels trustworthy.

% [GUIDELINE 2026-07-25 | testing (15%) | ACTION] The second half of this subsection is currently
% lost inside the comment block below: the within-subject user study, and with it the ONLY reference
% anywhere in the dissertation to Appendix~\ref{chap:ethics-statement}. As it stands the ethics
% appendix is orphaned -- it renders in the PDF but no sentence in the body points to it, which
% reads as an oversight rather than a considered position. The [FLAG] note below already records
% this. ACTION: reinstate the second study and the ethics reference as live text (the sentences
% beginning "The second validates the system: a within-subject study..." through "...neither of
% which the present work needed because it collected no human data (Appendix~\ref{chap:ethics-statement})").
% Keep Krippendorff's alpha and Gwet's AC1 commented out until you are comfortable explaining them,
% but the sentence can still say "chance-corrected agreement statistics" without naming them, as it
% already does. Consider a second pointer to the ethics appendix from Section~\ref{sec:limitations}.
The human work divides into two affordable studies. The first validates the instrument based on a modest set of transcripts graded independently by humans from different backgrounds and perspectives, which are compared against the judge's verdicts per principle using chance-corrected agreement statistics, turning the judge's accuracy from a design argument into a reported number.
% [FLAG] Your comment below commented out not only the two statistics but also the second (within-subject) user study and the ethics-appendix reference, so those are currently missing from the live text -- reinstate them if that was not intended (the reference to Appendix~\ref{chap:ethics-statement} was here).
% [EXPLAIN -- the two statistics you set aside, in plain terms] Both answer one question: how
% much do two independent graders agree, once the agreement you would expect from pure chance is
% removed? A raw ``they agreed 90\% of the time'' is misleading when almost everything gets the
% same label, because graders would agree that often by luck alone.
%   - Krippendorff's alpha: the general-purpose agreement coefficient. 1 means perfect agreement,
%     0 means chance-level, below 0 means worse than chance. It copes with any number of graders,
%     missing grades, and different scale types, which is why it is the usual default.
%   - Gwet's AC1: an agreement coefficient built to avoid the well-known ``paradox'' where alpha
%     (and Cohen's kappa) collapse towards 0 when nearly every item gets the same label -- which is
%     exactly this study's situation, since most sub-1B answers fail. AC1 stays stable there, so it
%     is reported beside alpha as the more trustworthy number for this case.
% If you reinstate them, keep the plain gloss in the live sentence and cite the two only in parentheses.
% (need to understand them first) Krippendorff's $\alpha$ \cite{krippendorff2004contentanalysis} and Gwet's AC1 \cite{gwet2008interrater}, the measurement would turn the judge's accuracy from a design argument into a reported number. The second validates the system: a within-subject study on the order of twenty participants, each carrying out the same personal-assistance tasks with two conditions and rating each on trust and satisfaction. The within-subject design, in which every participant sees both systems, controls for the wide differences between individuals, so even a small sample can detect whether the behavioural gains this work measured actually register as trustworthiness to the people they are meant to serve. Either study would require ethical review and a participant data-handling protocol before recruitment, neither of which the present work needed because it collected no human data (Appendix~\ref{chap:ethics-statement}).

\subsection{GRPO and Process-Based Rewards}
\label{subsec:conclusion-future-grpo}

Reinforcement learning was set aside in this work by choice rather than by refutation, a scope decision taken with the supervisor rather than a conclusion drawn from evidence. It is therefore an open extension: rewarding correct behaviour directly, rather than only imitating it. The whole dissertation rests on supervised fine-tuning, which can only imitate the behaviour its examples demonstrate; reinforcement-learning methods such as Shao et al.'s GRPO and its relatives such as Yu et al.'s DAPO can in principle reward good \emph{outcomes} the examples never showed, exploring beyond the demonstrated distribution \cite{shao2024deepseekmathpushinglimitsmathematical, yu2025dapoopensourcellmreinforcement}.

% [APPLIED 2026-07-25] The Rainone misattribution was corrected in the live text at all six
% sites (introduction x2, background, state-of-the-art, conclusion, discussion x2). RAGEN
% \cite{wang2025ragenunderstandingselfevolutionllm} now carries the 0.5B RL-collapse evidence and
% Rainone is cited only for what it shows. The instruction block below is retained as the record of
% what was wrong and why; do NOT re-apply it.
% [GUIDELINE 2026-07-25 | background/LR (15%) | DROP-IN] CITATION MISATTRIBUTION -- site 4 of 4.
% In the live paragraph starting "The interesting question, and the one":
%   REPLACE  "Rainone et al.\ report, these methods are empirically unstable at 0.6B, where the
%             small model's noisy outputs make the reward signal unreliable
%             \cite{rainone2025replacingthinkingtoolusage}"
%   ->       "these methods are empirically unstable below one billion parameters, where the model's
%             noisy outputs make the reward signal unreliable: RAGEN reports multi-turn
%             reinforcement-learning collapse on a 0.5B model
%             \cite{wang2025ragenunderstandingselfevolutionllm}, and Beyond ReAct excludes
%             Qwen3-0.6B from its reinforcement-learning phase for the same reason
%             \cite{wei2025reactplannercentricframeworkcomplex}"
% This also resolves the run-on the POLISH block above flags, so the two edits can be applied as one.
% Once all four sites are corrected, \cite{rainone2025replacingthinkingtoolusage} should remain only
% where it supports small-model reasoning brittleness -- principally the fourth hypothesis.
The interesting question, and the one the safety-tax finding sharpens, is whether reinforcement learning would \emph{sidestep} the trade-off or merely \emph{relocate} it: would rewarding outcomes recover the reasoning that imitation crowds out, or would the same fixed budget reassert itself under a different objective? This cannot be said, because it was not tested: these methods are empirically unstable below one billion parameters, where the model's noisy outputs make the reward signal unreliable, RAGEN reporting multi-turn reinforcement-learning collapse on a 0.5B model \cite{wang2025ragenunderstandingselfevolutionllm} and Beyond ReAct excluding Qwen3-0.6B from its reinforcement-learning phase for the same reason \cite{wei2025reactplannercentricframeworkcomplex}. As small-model reinforcement learning matures and stabilises, putting GRPO or a process-reward variant on the same bench is the natural next thing to try.

% ==================================================================================================
% [DRAFT 2026-07-25 | testing/conclusions (15%) -- STRONGEST FUTURE-WORK ADDITION AVAILABLE]
% GAP: Section~\ref{subsec:conclusion-future-user-studies} (validate the judge against humans) and
% this GRPO subsection are currently two unconnected loose ends. They are in fact one programme, and
% saying so converts Future Work from a list of things not done into a designed research trajectory
% -- which is exactly what the 70+ band means by "limitations/future research areas CLEARLY DEFINED".
% It also retro-justifies the whole judge apparatus: the instrument built to MEASURE becomes the
% instrument that TRAINS, so the effort spent on rubrics, anchors and bias controls pays twice.
% Note the ORDER matters and the draft makes the danger of getting it wrong explicit, which is the
% critical-analysis move rather than the enthusiastic one.
% Place as the closing passage of this subsection, or promote to its own subsection titled
% "From Measurement to Training Signal". To apply, delete the "% " prefixes.
%
% These two lines of future work meet, and stating the connection turns two loose ends into a single programme. The obstacle to reinforcement learning at this scale is not only the instability of the algorithms but the absence of a reward signal that is cheap, repeatable and aligned with the constitution: a person cannot sit in the loop for every sampled rollout, and a hand-written rule is too brittle to score behaviour rather than form. The judge built for this dissertation is already most of such a signal. It scores behaviour rather than wording, it grades against written per-principle criteria with anchored endpoints and a worked exemplar, it returns the same verdict on the same input often enough to be averaged, and every verdict is persisted with its evidence. What it lacks is validation: until its agreement with human markers has been measured, a reward derived from it would optimise the judge's preferences rather than a person's.
%
% The order of work therefore follows from the gap rather than from convenience. The judge is validated against human graders first (Section~\ref{subsec:conclusion-future-user-studies}), and only a judge that has cleared that bar is used as the reward model for reinforcement learning. Carried out in that order, the measurement apparatus of this dissertation becomes the training signal of the next: the same written constitution that defines the probes would define the reward, and the same rubrics that grade an answer would score a rollout, so that what the model is trained towards and what it is tested against derive from one document a person can read and contest. Carried out in the other order, an unvalidated judge would quietly become the operative definition of trustworthy behaviour, which is precisely the substitution the scrutability argument of Section~\ref{sec:constitutional-ai} exists to prevent, and the resulting system would be optimised against a standard nobody had ever checked.
% ==================================================================================================


\subsection{A Fast Interaction Layer over the On-Device Reasoner}
\label{subsec:conclusion-future-interaction}

The Thinker--Executor split of this work divides labour by \emph{role}, reasoning against acting, and runs the two sequentially, so the user waits for the whole chain before seeing anything; Chapter~\ref{chap:discussion} measured that hand-off latency as the dual model's main cost. A different split, by \emph{tempo}, has since been demonstrated at the frontier. Thinking Machines Lab's TML-Interaction-Small pairs a fast interaction model that keeps a conversation flowing in two-hundred-millisecond micro-turns with an asynchronous background model that handles the slow reasoning, web browsing and tool calls, the two sharing a single context so the interaction model can weave the background result into the live reply once it is ready \cite{thinkingmachines2026interaction}.

That fast-and-slow decomposition is a natural next axis for the architecture here. The on-device Thinker could run as the deliberate background reasoner while a lightweight interaction layer keeps the user engaged, hiding exactly the hand-off latency this work measured. The open question is  whether a fast-and-slow interaction split can be sustained with both halves entirely on the device, without the asynchronous reasoner quietly becoming a remote one.

\section{Closing Remarks}
\label{sec:conclusion-closing}

This dissertation asked whether a small language model, running entirely on a user's device, could be made trustworthy and personalised without sending private data to the cloud. The answer is a qualified one: the goal is partly within reach, and the contribution is a principled means of stating precisely which part, for any given model. The thread running through the work is that trustworthiness cannot be assumed in such models --- a model trained only to predict the next token states falsehoods in the same confident register as facts, cannot be relied upon to compute, and may memorise and later leak what it was shown. This work names those flaws as principles and measures how faithfully a small on-device model holds to them. It finds that safety and structure can be taught, that grounding buys what training cannot, and that visible reasoning is displaced as the price; the boundary between what is reachable and what is not is itself the result. The question was never whether a small model can comply, but where it stops complying, and whether that boundary can be measured, explained and designed against. As small models improve, that boundary is likely to move; the method for locating it, and the safety tax it exposes, should prove more durable.
